El proyecto tiene como principal alcance el desarrollo de un sistema de apoyo a la decisión orientado a la predicción de la demanda de los productos de una farmacia hospitalaria, esto utilizando los datos históricos de consumo como principal fuente de información. En el caso de estudio, el catálogo incluye medicamentos, dispositivos médicos, productos de nutrición y medicamentos controlados. El sistema permitirá estimar la cantidad mensual requerida de cada producto mediante la implementación de modelos de predicción de demanda, integrando funcionalidades como visualización de resultados, generación de reportes y apoyo en la toma de decisiones relacionadas con la planificación del abastecimiento. Asimismo, el sistema proporcionará al usuario cálculos como el punto de reposición, con el fin de contribuir a una mejor organización y control de los recursos disponibles en los entornos hospitalarios.

El desarrollo del proyecto presenta algunas restricciones que deben ser consideradas. En primer lugar, el sistema se basará en un conjunto de datos históricos reales de la Clínica Nuestra Señora de los Remedios que abarca solo 12 meses (de julio de 2021 a junio de 2022). Esto quiere decir que el modelo no cuenta con un año completo de historial adicional para comparar el mismo mes entre dos años distintos, lo que puede afectar su capacidad para capturar la estacionalidad, los comportamientos actuales o los cambios en la demanda a lo largo del tiempo.